\documentclass[12pt]{article}
\usepackage[utf8]{inputenc}
\usepackage[T1]{fontenc}
\usepackage[margin=1in]{geometry}
\usepackage{setspace}
\usepackage{xcolor}
\usepackage[colorlinks=true, urlcolor=blue!60!black, linkcolor=blue!60!black]{hyperref}
\setlength{\parindent}{0em}
\setlength{\parskip}{0.6em}
\renewcommand{\baselinestretch}{1.0}
\begin{document}
\begin{center}
{\LARGE \textit{NLRB v. Noel Canning}, 573 U.S. 513 (2014)\par}
\vspace{0.5em}
{\large Case Brief --- Two-Page Edition\par}
\end{center}

\section{Memory Jogger}
Recess appointments may be made in long intra-session recesses and for pre-existing vacancies, but the Senate is in session when it says it is, so the 2012 NLRB appointments were invalid (\href{https://www.law.cornell.edu/supremecourt/text/12-1281}{Cornell LII}).

\section{Detailed Case Facts}
The NLRB ruled against Noel Canning, a Pepsi-Cola bottler, which argued that the Board lacked a quorum because three members had been invalidly appointed. The Senate had scheduled ``pro forma session[s]'' every Tuesday and Friday over the holidays, with no business to be transacted; on January 4, 2012, during a three-day gap, President Obama recess-appointed the three members.

\section{Procedural History}
The D.C. Circuit vacated the Board's order, limiting the Clause to inter-session recesses and to vacancies arising during them. The Supreme Court affirmed on narrower grounds in 2014.

\section{Judicial Votes}
9--0 in judgment, 5--4 on the Clause's scope. Breyer wrote for Kennedy, Ginsburg, Sotomayor, and Kagan; Scalia concurred in the judgment, joined by Roberts, Thomas, and Alito.

\section{Holding}
The Clause covers intra-session recesses and pre-existing vacancies, but a recess of three days or fewer is too short and one under 10 days is presumptively too short. Pro forma sessions count because the Senate retained the capacity to act, so the appointments were invalid.

\section{Analysis of Opinions}
\textbf{Majority:} Where the text is ambiguous, long historical practice deserves ``great weight.'' Founding usage and two centuries of practice support intra-session appointments and pre-existing vacancies; the Adjournments Clause and practice support the 3- and 10-day lines. Courts must accept the Senate's own record, and the Senate actually passed a law during one pro forma session. \textbf{Scalia:} the text means only inter-session recesses and vacancies arising then; acquiescence cannot alter structural rules, and the majority's ``adverse possession'' theory favors the President.

\section{Examples: Future Applications}
\textbf{Same-Side:} an appointment during a three-day gap between pro forma sessions is invalid. \textbf{Opposite-Side:} an appointment during an 18-day summer adjournment to fill an earlier vacancy is valid. \textbf{Fence-Sitter:} an appointment after a catastrophe during a nine-day recess.

\section{Critique}
Breyer is functionalist on recesses and vacancies but formalist on pro forma sessions. Scalia's warning that treating Senate silence as acquiescence ratchets toward executive power remains potent, and the 10-day line appears nowhere in the text.

\section{Key Quotations}
\begin{quote}
``We hold that, for purposes of the Recess Appointments Clause, the Senate is in session when it says it is, provided that, under its own rules, it retains the capacity to transact Senate business.''
\end{quote}
\noindent --- Justice Breyer, for the Court, 573 U.S. at 550.

\end{document}
